\section{Finite support factors in the MPU--QCA index comparison}

\begin{lemma}[Least left support of the actual two-link output family]
    \label{lem:mpu_actual_two_link_left_support}
    \lean{Matrix.IsUnitaryBetween.conjStarAlgEquiv,
        Matrix.IsUnitaryBetween.conjStarAlgEquiv_apply,
        MPOTensor.twoLinkMiddleLift,
        MPOTensor.twoLinkOutputMap,
        MPOTensor.twoLinkOutputMap_apply,
        MPOTensor.twoLinkLeftOutputMap,
        MPOTensor.twoLinkRightOutputMap,
        MPOTensor.twoLinkOutputMap_inverse_conj_kronecker,
        MPOTensor.leftSupportAlgebra_twoLinkOutputMap,
        MPOTensor.twoLinkLeftSupportEquiv,
        MPOTensor.TwoSiteStandardFormData.left_div_phys_eq_phys_div_right}
    \notready
    Let $d,r,\ell$ be positive and let
    $u:\mathbb C^d\otimes\mathbb C^d\to\mathbb C^\ell\otimes\mathbb C^r$
    and $v:\mathbb C^r\otimes\mathbb C^\ell\to\mathbb C^d\otimes\mathbb C^d$
    be unitary. The actual output family
    \[
        A\longmapsto(v\otimes v)
        (I_r\otimes uAu^\dagger\otimes I_\ell)(v\otimes v)^\dagger
    \]
    is a star-homomorphism. Its least left support algebra is
    $v(I_r\otimes M_\ell)v^\dagger$, and hence has matrix size $\ell$.
    This is the finite output-family step in
    \cite[Appendix, lines~2322--2334]{Cirac2017MPU}.
    The adjoint order follows the main-text direction of $v$; the
    Appendix correction and its distinct logarithmic factor-of-two
    correction are recorded in \cite{gap:mpu_qca_index_normalization}.
\end{lemma}
\begin{proof}
    Change the arbitrary input by the inverse of $u$. An intermediate
    product $A_\ell\otimes B_r$ then gives exactly
    $v(I_r\otimes A_\ell)v^\dagger\otimes
        v(B_r\otimes I_\ell)v^\dagger$. Matrix units span the intermediate
    algebra, so every left coefficient belongs to the first factor.
    Conversely choose $B_r=I_r$; a diagonal right coefficient recovers
    every element of the first factor. Leastness gives equality.
    Conjugation by $v$ and the faithful identity-tensor embedding
    identify it with $M_\ell$.
\end{proof}

\begin{lemma}[Coordinate transport of least finite left support]
    \label{lem:mpu_qca_left_support_coordinate_transport}
    \lean{Matrix.bipartiteSlice_reindex_prod,
        Matrix.leftSupportAlgebra_map_matrixReindexStarAlgEquiv}
    \notready
    Independent coordinate bijections on the two factors transport the
    least left support algebra by the first-factor bijection.
\end{lemma}
\begin{proof}
    A left coefficient in the new coordinates is exactly the
    first-factor coordinate transport of the old coefficient, with
    its two right indices relabelled. Apply the least-support
    characterization in both directions.
\end{proof}

\begin{lemma}[Intersection of finite quasi-local supports]
    \label{lem:mpu_qca_support_intersection}
    \lean{SpinChain.commute_map_localObservable_of_singletons,
        SpinChain.QuasiLocalSupportedIn.inter}
    \notready
    An observable supported on two finite regions $\Lambda$ and $\Gamma$
    is supported on $\Lambda\cap\Gamma$.
\end{lemma}
\begin{proof}
    Any site outside the intersection lies outside at least one
    of the two supports. The observable therefore commutes with
    its one-site algebra. Finite one-site generation gives commutation
    with every algebra disjoint from the intersection, and the
    outside-commutant characterization recovers exact support.
\end{proof}

\begin{lemma}[Actual complete-interval action and its four-block output]
    \label{lem:mpu_actual_complete_interval_four_output}
    \lean{SpinChain.quasiLocalIntervalObservable_chainWindowOperator,
        SpinChain.bufferedMPUQuasiLocalAut_central_interval,
        SpinChain.bufferedMPUQuasiLocalAut_central_interval_four_output}
    \notready
    \uses{lem:mpu_supplied_standard_form_actual_qca}
    Under the supplied positive-block standard-form hypotheses, the
    actual completed action sends every matrix on the complete
    original-site interval $[a+2k,a+4k)$ to the two-link output matrix
    on $[a+k,a+5k)$, with the indicated block-coordinate identifications.
    In particular, its output support is this four-block interval.
\end{lemma}
\begin{proof}
    Choose the actual ambient chain $[a,a+8k)$ so that every input
    site has the required buffers. The finite-chain evaluation of
    the completed action and the actual standard-form identity give
    the fixed six-block output. Remove the two identity exterior
    blocks using its exact four-block placement and canonical
    interval inclusions.
\end{proof}
